\chapter{Base histórica e atualização teórica}\label{o-que-o-tg-oferece-e-o-que-precisa-ser-atualizado}

O trabalho de graduação de Wayne Leonardo Silva de Paula
\autocite{tg2003} constitui a base histórica. O núcleo aproveitável
encontra-se nos capítulos 4 a 6 e no apêndice Maple:

\begin{itemize}
\tightlist
\item
  \textbf{Capítulo 4, páginas impressas 22--25:} equações linearizadas
  da teoria de Visser, conservação e propagação dispersiva,
  especialmente as equações (4.14)--(4.22).
\item
  \textbf{Capítulo 5, páginas 27--38:} formalismo de tetradas,
  aproximação quase nula, classificação E(2), equações (5.58)--(5.61) e
  matriz das componentes elétricas do tensor de Riemann.
\item
  \textbf{Capítulo 6, páginas 39--42:} comparação entre faixas de
  frequência e sugestão de maior relevância de efeitos massivos em
  baixas frequências.
\item
  \textbf{Apêndice A, páginas 43--51:} procedimento computacional que
  pode ser reconstruído e verificado em ferramentas atuais.
\end{itemize}

O resultado central já foi publicado por de Paula, Miranda e Marinho em
\emph{Classical and Quantum Gravity} \textbf{21}, 4595--4606 (2004).
Assim, reproduzir as seis polarizações ou simplesmente atualizar os
números do TG constitui uma etapa de formação e reprodução, mas não a
contribuição original de um novo artigo \autocite{depaula2004}.

Há quatro distinções fundamentais para o mestrado:

\begin{enumerate}
\def\labelenumi{\arabic{enumi}.}
\tightlist
\item
  \textbf{Onda massiva e aproximação quase nula.} Usar uma tetrada nula
  é permitido; o problema é aplicar relações simplificadas que
  pressupõem propagação nula fora de seu domínio. Hyun, Kim e Lee já
  desenvolveram expressões exatas para polarizações nulas e não nulas e
  discutiram limitações de trabalhos anteriores, incluindo a referência
  de 2004. Uma nova derivação, isoladamente, não seria inédita \autocite{hyun2019}.
\item
  \textbf{Deformações observáveis e graus de liberdade.} A matriz
  simétrica \(R_{0i0j}\) possui seis componentes, mas os vínculos de uma
  teoria podem relacioná-las. Um campo de Fierz--Pauli massivo em quatro
  dimensões tem cinco graus de liberdade; sua helicidade zero pode
  produzir deformações transversal e longitudinal relacionadas. Isso
  difere de assumir seis amplitudes independentes \autocite{liang2021}.
\item
  \textbf{Consistência da teoria.} O termo massivo não Fierz--Pauli da
  formulação histórica exige atenção ao modo escalar com sinal cinético
  inadequado, o \emph{ghost}. A contagem geométrica de polarizações não
  demonstra estabilidade. O modelo de Visser será uma referência
  histórica; o modelo físico de comparação será o setor linear de
  Fierz--Pauli, com discussão de seu domínio e de sua relação com
  construções modernas. A estabilidade não linear de dRGT ou
  Hassan--Rosen não pode ser transferida automaticamente para Visser \autocite{visser1998}, \autocite{park2010}, \autocite{derham2011},
  \autocite{hassan2012}.
\item
  \textbf{Amplitude de curvatura e detectabilidade.} Comparar
  coeficientes que multiplicam perturbações métricas não substitui
  calcular excitação pela fonte, resposta instrumental e ruído. A
  hipótese de componentes métricas com amplitudes semelhantes,
  explicitada no artigo de 2004, não é uma previsão geral sobre fontes
  astrofísicas. A classificação E(2), o tipo algébrico de Petrov e o
  número de graus de liberdade também devem ser apresentados como
  conceitos distintos.
\end{enumerate}

Uma auditoria inicial deverá fixar convenções de sinal, normalização das
tetradas, distinção entre \(h_{\mu\nu}\) e sua versão de traço
invertido, unidades e a definição de massa nas equações (4.17)--(4.22).
As tabelas do capítulo 6 devem ser recalculadas com unidades explícitas.
Não é necessário presumir que todas as diferenças sejam erros físicos;
algumas podem decorrer de notação ou normalização.
